% Ad Quintum Fratrem 3.8 --- headnote
% ad-quintum-fratrem-03-08, mid- to late November 54 BC, Rome

\textit{Cicero to Quintus, written from Rome in mid- to
late November 54 BC. The Perseus dateline reads only
\textit{Scr. Romae ex. m. Nov. a. 700 (54)}, end of November;
the body fixes the writing more precisely: \textsection 5
records the funeral of Serranus's son on 24 November as a
recent event, and \textsection 6's news about the games and
Milo's preparations places the letter in the immediate
sequence after Q. fr. 3.3 and 3.4 but before the year's
close.}

\textit{Two letters from Quintus had reached Rome: the earlier
one ``full of vexation and complaints'' (and another given to
Labienus the day before), the more recent of which had set
those distresses aside. \textsection 1 is the brother-counsel
on bearing what life on Caesar's staff in Gaul is asking of
him: the design of the joining had not been the small
advantages but ``the most secure protection out of the
goodwill of the best and most powerful man, towards the whole
standing of our dignity.'' \textsection 2 adds the caution
about what should not be entrusted to letters, and about the
courier route --- the Nervii are in the field, the new
winter-camp positions of the legions are not yet known to
Cicero, and the letters need to go through Caesar's couriers
or Labienus's. The siege of Quintus's camp at Atuatuca by
Ambiorix would come within the next several weeks.}

\textit{\textsection 3 contains the two principal political
moments of the letter: Caesar's ``courage and gravity in his
deepest grief'' (the bereavement of Julia, of which Cicero
had received Caesar's letter at the moment of Q. fr. 3.1
\textsection 17), and the news that the candidates of the
year have been ``relieved of their vexation'' --- the
canvassing-prosecutions of late October are weakening, and
Messala and Domitius are now counted as the probable consuls
of 53. The poem to Caesar that Cicero had abandoned at
Q. fr. 3.1 \textsection 11 will now be taken up again, on
the leisure days of the November thanksgivings (the
\textit{supplicationes} voted to Caesar for the British
expedition), and given to Caesar through Quintus.
\textsection 4 is the running report on the dictatorship
talk: the elections postponed, the interregnum coming;
Pompey openly denies that he wants it, though to Cicero in
person he had not denied it; C. Lucilius Hirrus is the
tribune working to propose the law, and his self-regard
(``without a rival,'' a Catonian tag) is mocked. \textsection
6 is the report on Milo, who will give the ruinously
extravagant games of 54 in spite of Cicero's advice not to:
Pompey has thrown him over for Cotta, and if Pompey is made
dictator Milo ``almost gives himself up for lost'' --- the
shape of things that, with the killing of Clodius two months
later, will produce the Pro Milone trial of April 52.}

\textit{P. Crassus's death at Carrhae had become known in
Rome at the end of October (the first piecemeal reports), and
the heavy public mood the year closes in --- ``the rumour of
a dictatorship is unpleasant to the good men, and even more
so to me are the things they are saying'' (\textsection 4) ---
is the cumulative effect of the Carrhae news, the broken
elections, and the dispersal of the legions for the worst
Gallic winter of the campaign. The final domestic note
(\textsection 5), the funeral oration delivered by Serranus
on a script Cicero had written for him, is the unobtrusive
private favour that runs alongside the political letter.}
